\begin{proposicion}[Reconstrucción de la clase interior]\label{meso:reconstruccion}
Existen constantes $C_{\rm int},\eta_{\rm int},\gamma_{\rm int}>0$, que dependen sólo del perfil, con
la propiedad siguiente.
Supongamos que los observables de esta sección, es decir, las proporciones por intervalos, las fibras
y los dieciséis términos del rectángulo, están dentro de sus tolerancias y que
$D_{\max}(k)<\alpha_{\rm geo} k$ con $\alpha_{\rm geo}=\varepsilon/(4\|W_0\|_\infty)$. Si el paso
$a_{\rm mes}$, la tolerancia
$\varrho_{\rm mes}$ y $\varepsilon$ son suficientemente pequeños y $k$ es grande, el punto de
Bethe del tablero existe, cumple con desigualdad estricta las cotas de difusión, vacancia y brecha
con esas constantes, y la masa de cada diagonal es menor que $\varepsilon$. En las casillas
disponibles $p_e/q_e=1+O(\varrho_{\rm mes}+\varepsilon+1/k)$, con constante dependiente del perfil. Además,
para esos $\varepsilon$ y $k$ grande, todo hijo de un tablero cuyo punto de Bethe cumple esas cotas con
masa diagonal a lo
sumo $\varepsilon$ cumple las cotas exteriores, y la masa de cada una de sus diagonales es a lo sumo
$\varepsilon_{\rm ext}=2\varepsilon$.
\end{proposicion}
\begin{proof}
La suma de cada fila y columna de $q$ es $1+O(\varrho_{\rm mes})$ por las fibras. Para $V$ se aproxima
$\nu_1(r,\cdot)$, o $\nu_1(\cdot,c)$, por una constante en cada intervalo. La normalización continua
del perfil, su regularidad de Lipschitz y la suma de las tolerancias de los intervalos dan grados
$1+O(a_{\rm mes}+\varrho_{\rm mes})$ uniformemente en la coordenada de la línea distinguida. Fijemos
${0<\theta<\min(c_\nu,1)/2}$.
La matriz $V-\theta\mathbf1\mathbf1^T/k$ es no negativa y tiene sumas de filas y columnas
${\le1-\theta+O(a_{\rm mes}+\varrho_{\rm mes})}$. Su norma de operador está acotada por esa misma
cantidad, por la
desigualdad de
Schur. La aproximación de rango uno y la cota para $q-V$ dan
\[
 \sigma_2(q)\le1-\theta+O(a_{\rm mes}+\varrho_{\rm mes})+\|q-V\|_{\rm op}.
\]
Si $D_R,D_C$ son las
matrices diagonales de
los grados de $q$, la matriz $D_R^{-1/2}qD_C^{-1/2}$ difiere de $q$ en norma de operador en
$O(\varrho_{\rm mes})$. Con las tolerancias anteriores suficientemente pequeñas, su segundo valor
singular es $\le1-\theta/2$. Denotemos por $L_q$ el laplaciano ponderado de filas y columnas y
pongamos $D=\operatorname{diag}(D_R,D_C)$ y $d_{\min}=\min_v D_{vv}$. Los bloques no diagonales de
la matriz de adyacencia normalizada son $D_R^{-1/2}qD_C^{-1/2}$ y su traspuesta, de modo que sus
valores propios son los valores singulares con ambos signos. Su mayor valor propio es uno, con vector
propio $D^{1/2}\mathbf1$. Así, el laplaciano normalizado $D^{-1/2}L_qD^{-1/2}$ tiene brecha
$\lambda=1-\sigma_2(D_R^{-1/2}qD_C^{-1/2})\ge\theta/2$. Para todo vector real $f$, al restar la
media ponderada por los grados, la desigualdad de Rayleigh da
\[
 f^TL_q f\ \ge\ \lambda\min_{a\in\mathbb R}\sum_v D_{vv}(f_v-a)^2
 \ \ge\ \lambda d_{\min}\min_{a\in\mathbb R}\sum_v(f_v-a)^2 .
\]
Para $f\perp\mathbf1$ el último mínimo es $\|f\|_2^2$, por lo que
la brecha de $L_q$ es al menos $d_{\min}\lambda$. Aquí todos los grados son $1+O(\varrho_{\rm mes})$.

Detallamos cómo pasar de estos observables al punto de Bethe. En el soporte de $q$, sea $B$ la matriz
de incidencia de filas y columnas de la Sección~\ref{macro:reconstruccion} y
$P_q=\operatorname{diag}(q_e)$. Para el reequilibrado resolvemos
$(BP_qB^T)z=\mathbf1-Bq$ en el complemento ortogonal del gauge $(\mathbf1,-\mathbf1)$. El segundo
miembro es ortogonal a ese vector. El cambio de signo en las coordenadas de columna convierte la
matriz en el laplaciano de filas y columnas de pesos $q$. La brecha anterior y los grados acotados
superior e inferiormente dan $\|z\|_2/\sqrt{k}=O(\varrho_{\rm mes})$. En la ecuación de cada vértice,
el término
de vecinos es a lo sumo $(\sum_j q_{ij}^2)^{1/2}\|z\|_2=O(\varrho_{\rm mes})$, pues
$q_{ij}\le\|W_0\|_\infty/k$;
análogamente en columnas. Al dividir por el grado resulta $\|z\|_\infty=O(\varrho_{\rm mes})$. Por tanto
\[
 \bar q=q+P_qB^Tz,\qquad B\bar q=\mathbf1,\qquad
 \bar q_e/q_e=1+O(\varrho_{\rm mes}).
\]
Las cotas de difusión y de brecha siguen cumpliéndose con desigualdad estricta por comparación de los
pesos. Como $D_{\max}(k)<\alpha_{\rm geo} k$ con $\alpha_{\rm geo}=\varepsilon/(4\|W_0\|_\infty)$, cada
diagonal de las
dos familias tiene masa de $q$
a lo sumo $\varepsilon/4$ y masa de $\bar q$ a lo sumo $\varepsilon/2$ tomando $\varrho_{\rm mes}$
suficientemente pequeño con independencia de $\varepsilon$.

Para precisar la inversión relativa, pongamos $P=\operatorname{diag}(\bar q)$ y
\[
 \mathcal T=\{x:BPx=0\},\qquad
 \Pi=I-B^T(BPB^T)^+BP ,
\]
donde el superíndice $+$ denota la inversa en el complemento ortogonal del gauge, prolongada por cero. Esta
proyección sobre $\mathcal T$ anula $\operatorname{im}B^T$, el espacio de los multiplicadores de fila
y de columna, y tiene norma de $\ell^\infty$
acotada uniformemente: se aplica la estimación anterior del laplaciano con segundo miembro $BPv$, de
norma $\le\|v\|_\infty$ porque $B\bar q=\mathbf1$. En las variables factibles ${p=\bar q\odot(1+x)}$,
$x\in\mathcal T$, la ecuación KKT es $F(x)=-\Pi\nabla H_B(p)=0$, donde $H_B$ es la entropía de
Bethe \eqref{bethe}, cuyo subíndice es el tablero y no la matriz de incidencia. Módulo multiplicadores
de fila y de columna, su gradiente tiene componentes
$-\log p_e-3\log(1-p_e)+\log h_{D(e)}+\log h_{A(e)}$. En $\bar q$ se tiene
$\log q\in\operatorname{im}B^T$
por $W_0=b_0 r_Rr_C$, y $\log(\bar q/q)=B^Tz+O(\varrho_{\rm mes}^2)$ en norma uniforme. Los términos de
vacancia
y de $1-\bar q_e$ son $O(\varepsilon+1/k)$. Por consiguiente $\|F(0)\|_\infty\le
C(\varrho_{\rm mes}^2+\varepsilon+1/k)$.

En la bola $\|x\|_\infty\le t<1/2$ la derivada de $\log p$ en variables relativas es la diagonal
$1/(1+x_e)$; su diferencia con la identidad es $O(t)$. La derivada de $\log(1-p_e)$ tiene norma
$O(1/k)$, y cada derivada de $\log h_f$ tiene norma $\le\sum_{e\in f}\bar q_e/h_f=O(\varepsilon)$. La
proyección es la identidad en $\mathcal T$, de modo que $\|DF(x)-I\|_{\infty,\mathcal T}\le
C(t+\varepsilon+1/k)$. Fijamos $t$ pequeño y después $\varrho_{\rm mes},\varepsilon,1/K_0$ suficientemente
pequeños
para que $x\mapsto x-F(x)$ sea una contracción con constante $\le1/2$ en la bola y
$2\|F(0)\|_\infty\le t$. El punto fijo satisface
\[
 \|p/\bar q-1\|_\infty\le 2\|F(0)\|_\infty,\qquad
 p_e/q_e=1+O(\varrho_{\rm mes}+\varepsilon+1/k).
\]
Todas las razones se toman en casillas disponibles. El punto es el único estacionario del dominio
convexo $0<p_e<1/4$, $h_f>0$, ya que la hessiana negativa de la entropía de Bethe tiene allí diagonal
positiva
$1/p_e-3/(1-p_e)$ más formas positivas semidefinidas de las líneas. Como $p_{\max}\le C/k\le e^{-4}$
para $k$ grande,
el cálculo escalar de la Proposición~\ref{macro:maximo} muestra además que el punto construido es
el único máximo global. Es, por tanto, la rama seleccionada en la representación de
\cite[Teorema~8.1 y el párrafo que lo sigue]{pred}. Obtenemos las cotas de difusión, brecha y
vacancia con desigualdad estricta, y masa diagonal $<\varepsilon$. La construcción usó solamente
observables del tablero alcanzado. Para cualquier hijo de un último padre regular, aun cuando falle
uno de esos observables, la herencia de \cite[Secciones~6.5 y~8.2--8.3]{pred}, que es el punto (I1)
de la Sección~\ref{leyes}, da en las casillas que siguen disponibles una variación relativa
$1+O(1/k)$ respecto del padre, con constantes de la clase exterior. Su masa diagonal es a lo sumo
$(1+O(1/k))\varepsilon\le\varepsilon_{\rm ext}=2\varepsilon$ tras aumentar $K_0$.
\end{proof}
